\section{Introduction}
Most of the microbial world has never been grown in the laboratory. What we know about it comes increasingly from genomes
assembled out of environmental DNA; the Genomes from Earth's Microbiomes (GEM) catalog alone holds \NGEMTotal of them
\citep{nayfach2021}. For these organisms the genome is the only record of their physiology, and one of the first things to ask
of it is how fast the organism can grow. The maximal growth rate separates microbes that bloom on a pulse of resources from
those that persist at low activity, and it is needed to interpret abundance and replication estimates from metagenomes
\citep{vieirasilva2010,roller2016,kirchman2016}.

The genome does carry a signal. Fast growers concentrate the codon usage of their most highly expressed genes, above all the
ribosomal proteins, on a few preferred codons, and the strength of that bias tracks the maximal growth rate
\citep{vieirasilva2010}. gRodon turns codon usage statistics into a predicted doubling time \citep{weissman2021} and has been
applied to whole metagenomes \citep{weissman2022}. Phydon adds the doubling times of close relatives on the GTDB tree
\citep{xu2025,parks2022}. gRodon was trained on the compilation of \citet{madin2020}, which pools values from many studies and
databases.

Whether such predictions can be trusted for an uncultured genome depends on three things that are easy to overlook. The first
is distance: \GEMShareNoGenus\% of high-quality GEM genomes have no labeled relative in their genus and \GEMShareNewOrder\% none
in their order, so what matters is accuracy at that distance from the training data, and the error of each individual
prediction. The second is temperature. Both tools take the measured growth temperature as an input, and temperature explains
much of the variation in growth rate, but a metagenome-assembled genome does not come with one. The third is the labels. In a
compilation, a doubling time measured far from optimal conditions looks just like a maximal rate, and an error shared by every
value from one study can pass for biology.

Here we predict the maximal doubling time, the growth temperature and a doubling time standardized to \TRefC\,\textdegree C
from the genome alone, and ask how far each prediction can be trusted for an uncultured genome
(Figure~\ref{fig:overview}). Our contributions are:
\begin{itemize}
\item An audit of the training labels (Algorithm~\ref{alg:audit}) that removed \NExcluded values that are not maximal growth
rates; \NExcludedKevbrinaWord of them, from one study of growth at low temperature, alone created an apparent
\KevFoldWith-fold effect of methane and ammonia oxidation on doubling time with relatedness controlled.
\item A nested evaluation with whole species, genera, families or orders held out (Algorithm~\ref{alg:nested}), which
attaches the error and an 80\% range of its novelty level to every prediction, with bootstrap intervals, a taxonomy
baseline and a shuffled-label control.
\item A comparison with the released gRodon and Phydon on the same genes and folds, with species and with orders held out,
tested per genome. For a new species our model without the temperature is as accurate as either tool given it
(\RelOursErr-fold against \RelGRodonErr- and \RelPhydonErr-fold); for a new order it trails gRodon given the temperature and
beats gRodon without it.
\item A test on \ExtN species described outside both training compilations, on which the model beats taxonomy and gRodon
but its ranges are too narrow.
\item A phylogenetically controlled test of whether slow growers are slow because they live on thin energy budgets
\citep{hoehler2013,lever2015}, which finds no support among cultured organisms.
\item A test on \MagNGenomes metagenome-assembled genomes of \MagNSpecies species, and predictions with errors for the
\NGEMPredicted high-quality GEM genomes.
\end{itemize}
